\documentclass[10.5pt]{article}
\usepackage[utf8]{inputenc}
\usepackage[T1]{fontenc}
\usepackage{geometry}
\geometry{
  paperwidth=5.5in,
  paperheight=8.5in,
  margin=0.65in,
  top=0.75in,
  bottom=0.75in
}
\usepackage{mathpazo}
\usepackage{microtype}
\usepackage{lettrine}
\usepackage{amssymb}
\usepackage{footmisc}

\usepackage{hyperref}
\hypersetup{
    colorlinks=true,
    linkcolor=black,
    citecolor=black,
    urlcolor=black
}

\pagestyle{plain}

\newcommand{\diamondrule}{%
  \begin{center}
    \vspace{-0.5em}
    \rule{0.25\textwidth}{0.4pt}\hspace{0.5em}\raisebox{-0.2ex}{\small$\blacklozenge$}\hspace{0.5em}\rule{0.25\textwidth}{0.4pt}
    \vspace{-0.5em}
  \end{center}
}

\begin{document}

% PAGE 1: TITLE PAGE
\begin{titlepage}
    \centering
    \vspace*{1.8in}
    {\scshape\footnotesize A STORY IN SIX SHORT SCENES\par}
    \vspace{0.6in}
    {\Huge\textbf{Loving People Well}\par}
    \vspace{0.3in}
    \diamondrule
    \vspace{0.4in}
    {\itshape\large Adam Matthew Steinberger\par}
    \vfill
\end{titlepage}

% PAGE 2
\newpage
\setcounter{page}{2}

\lettrine[lines=3,findent=3pt,nindent=0pt]{\textbf{M}}{ost} failures of love are not dramatic. They happen in ordinary hours, over dinner, in the car, on a Thursday, and they feel like nothing while they happen. Over months they do the damage. This is the story of one such hour, and of what slowly began to mend it.

\vspace{1.5em}
\diamondrule
\vspace{1.5em}

\begin{center}
    {\scshape\large The Phone at Dinner}\par
    \vspace{0.8em}
\end{center}

Nora loved Daniel, and would have said so. That was the trouble. She meant it, and still her phone lit up beside her plate while he explained, slowly, the way people do when it costs them, that the shop had lost its biggest client.

``Mm,'' she said, eyes on the screen. ``That's rough.''

Daniel nodded and talked about the weather.

Driving home she saw it the way you see a stain after the party. She had been fretting over a dozen small things when only one needed her, and he had been sitting right there.\footnote{Luke 10:40-42}

The next week she left the phone in the car. She asked about the shop, then the client, then what he was afraid of, and she waited out the silence the way you wait for a deep well to answer a dropped stone.\footnote{Prov 20:5} The fear was not the money. His father had built that shop, and Daniel could not bear to be the one who lost it. She would never have learned that by looking busy. So she began to learn his week, his work and his worries, because she was interested in him and not because she needed an audience.\footnote{Prov 27:23}

% PAGE 3
\newpage

\begin{center}
    {\scshape\large The Long Saturday}\par
    \vspace{0.8em}
\end{center}

Learning Daniel took longer than she liked. Some weeks he wanted company. Some weeks he wanted the door left open and nobody walking through it. She learned to walk at his pace, to go the second mile when he asked and not a step before,\footnote{Matt 5:41} and to know there is a time to speak and a time to sit quietly beside someone.\footnote{Eccl 3:7}

Then came a Saturday with no purpose at all. They played cards badly and laughed until a neighbor knocked. They traded old stories. Somewhere in the afternoon Nora told him about a failure from her twenties she had never told anyone, and the room went quiet, and then Daniel told her his. A cheerful heart does what medicine does, her grandmother used to say.\footnote{Prov 17:22} She had spent years believing love was labor. Rest, it turned out, was part of the work.\footnote{Matt 11:28; Mark 6:31}

Invite people into your story, and let them invite you into theirs.\footnote{1 Thess 2:8; Rom 12:15} Neither half works without the other.

\vspace{1.5em}

\begin{center}
    {\scshape\large Tuesday's Word on Friday}\par
    \vspace{0.8em}
\end{center}

The trouble with Tuesday was Friday. Daniel mentioned his daughter's surgery on Tuesday; by Friday Nora had invited him to dinner without a word about it. His face did not change. That was worse.

Forgetting, she decided, was a small abandonment. So she did something unglamorous. After each visit she spoke a voice memo into her phone: the money trouble, the surgery date, the joke he told. She dropped screenshots of his texts into one note. On Thursday nights a reminder told her what to ask on Friday. The phone that had stolen her attention at dinner now helped her keep it.

% PAGE 4
\newpage

On the next Friday she asked, ``How is Maddie after the surgery?''

Daniel stopped with his fork in the air. ``You remembered,'' he said, and his voice caught, because being remembered is a small way of being held.

She carried his name the way you carry something written on the palm of your own hand,\footnote{Isa 49:16} and was grateful, truly, every time she remembered him.\footnote{Phil 1:3}

\vspace{1.5em}

\begin{center}
    {\scshape\large Seventy Times}\par
    \vspace{0.8em}
\end{center}

Then Daniel snapped at her over nothing, and skipped her birthday dinner, and she caught herself keeping a ledger. Seven offenses by spring. She remembered an old exchange in which someone asked whether seven forgivenesses were enough and was told, not seven, but seventy times seven.\footnote{Matt 18:21-22} She tore the ledger up. She chose to believe the best of him first, and to hope, and to put up with a good deal.\footnote{1 Cor 13:7}

Mercy, though, had edges. When he asked to borrow money for the third time, she said no, gently. A friend's wounds can be trusted more than an enemy's kisses.\footnote{Prov 27:6} She told him the truth when nodding along would have been easier,\footnote{Eph 4:15} and she said yes only when she meant yes, so that when she promised to come Friday, she came.\footnote{Matt 5:37}

When Daniel finally saw a counselor, she cheered. She did not compete. People have a right to heal, and she never stood in the way.

Attention without memory, she decided, was a performance. Memory without mercy was a ledger. Mercy without edges was neglect. None of it made her hard, and none of it made her soft. It made her honest, which is the kinder thing.

% PAGE 5
\newpage

\begin{center}
    {\scshape\large Open Hands}\par
    \vspace{0.8em}
\end{center}

Control came last, and quietest. When Daniel drifted, Nora tightened her grip and drew up plans for his recovery, and the harder she held, the further he stepped back. So she set the whole weight down on someone larger than herself\footnote{1 Pet 5:7} and found her hands free enough to hold him lightly. She let him see her weakness. It turned out to be where the strength was.\footnote{2 Cor 12:9-10}

One evening he asked, half joking, whether he was a weight to her or a gift. She did not answer in words. The answer took a year of Tuesdays, because a tree is known by what it bears.\footnote{Matt 7:16-20} And she never held any of this over him. A framework like this binds only the one who adopts it. She tended her own eye long before she mentioned the splinter in his.\footnote{Matt 7:3-5}

\vspace{1.5em}

\begin{center}
    {\scshape\large Thirty Days}\par
    \vspace{0.8em}
\end{center}

I cannot prove any of this. Nobody measured Nora and Daniel, and the mess remains; it always will. But love, she concluded, is attention that stays, under mercy, held with open hands.\footnote{1 Cor 13:4-7}

Begin small, because small is where love lives, ordinary hour by ordinary hour. So try it. Choose one person. Put the phone away and learn one new thing about their week. Write it down. Remember it on Friday. Forgive one small thing. Then rest. After thirty days, ask whether they feel more known.\footnote{Ps 139:1-2}

Take your time. There is enough of it.\footnote{Ps 46:10; Matt 6:34}

\vspace{1.5em}
\diamondrule

\end{document}
